\subsection*{60. Connected topoi, geometric weak equivalences and initial axioms (pp.~148--151)}
In returning to test categories, Grothendieck distinguishes local connectedness, connectedness of the final object, and closure of connected objects under products. Local connectedness is expressed through a generating family of 0-connected objects; total 0-connectedness adds the product condition, while strict total 0-connectedness requires every nonempty object to have a section. The product test may be checked on generators. This applies the interval machinery of sections 54--59 to presheaf topoi, while the semisimplicial category without degeneracies provides a warning: its representables have no global sections, and on p.~149 Grothendieck deletes the proposed exception for the vertex. He asks, rather than proves, whether modelizing together with total 0-connectedness forces the stronger conditions. Under strict total asphericity he does prove that weak equivalences are geometric and that $h_W$ agrees with the canonical homotopy structure; the handwritten strengthening says that any connected generating family supplies the testing intervals, and the marginal reference points forward to p.~268. This equality reduces the modelizing question to the existence of an object with two distinct sections, or equivalently in the stated situation to connectedness of the Lawvere contractor. The topos of sets fails this condition, and p.~150 conjectures that it is the only counterexample under these hypotheses. The section then replaces cohomologically defined weak equivalences by a general saturated class $\underline W$ of functors, pulls it back through $i_A(F)=A_{/F}$, and lists 0-admissibility, invariance under product with the interval, and the final-object axiom. On p.~151 the final-object axiom is recognized as a consequence of the interval axiom; this shortens the initial list and leads directly to section 61's comparison theorem. A caution is needed for the intervening assertion that a category has a final object if and only if its identity admits an interval homotopy to a constant functor. The forward implication is valid, but the converse is false: the one-object category with endomorphism monoid $\{1,0\}$, where 0 is absorbing, has a natural transformation from its identity to its constant functor with component 0, yet its sole object is not final because it has two endomorphisms. Thus contractibility by such a transformation does not establish finality. This does not invalidate the deduction of the final-object axiom from interval invariance, which uses only the forward implication. The readers preserve the false converse, the literal superscript in $\Delta^1$, and the author's tentative questions; the guide separates their logical roles from the proved geometricity statement. The later argument in section 96 should also be read with the actual section-existence hypothesis: having only two subobjects of the final object alone does not imply that every nonempty object has a section.

\subsection*{61. The promised key result and a first list of axioms (pp.~151--156)}
Grothendieck resumes the comparison of a functor $i:A\to(Cat)$ with its induced $i^*:(Cat)\to\widehat A$, temporarily allowing two saturated arrow classes $\underline W$ and $\underline W'$. The desired equivalence of localizations is strengthened to the exact inverse-image condition $\underline W'=f_i^{-1}(\underline W)$, where $f_i=i_Ai^*$, so that successful sufficient conditions might also characterize the usual weak equivalences. The comparison diagram from pp.~96--98 is rebuilt using relative asphericity of $A$, contractibility of the categories $i(a)$, preservation of weak equivalences by fibered categories under condition e), and a slice criterion f). The decisive computation identifies the slice over $(a,p)$ of $A_{/C}\to A_{//C}$ with the category of elements of a pullback of $i^*(\underline{Fl}(C))\to i^*(C)$; thus the problem becomes universal membership of this last map in $W_A$. This develops the earlier homotopy and internal-Hom machinery and supplies the lemma reformulated in section 62. The text records its own weakening on p.~154: 0-admissibility b) is unnecessary for the intermediate comparison, since only $h_{(Cat)}\leq h_{\underline W}$ is needed. Two source slips require care. Diagram (11), p.~153, ends in $A$, although the intended comparison and relative-asphericity argument require the projection $A\times C\to C$, as the earlier diagram on p.~96 explicitly shows. The same paragraph changes the section reference to 57 while retaining pp.~96--98; readers should follow those pages for the diagram. Neither is repaired in the readers. Page 155's literal $f_Y$ and lower mark by $Y$ are retained, and the visible label $p$ on the pullback square, omitted by the witness, is restored. At p.~156 Grothendieck immediately withdraws the independent choice of the two arrow classes: the next section assumes $\underline W=\underline W'$ rather than claiming that this first axiomatization succeeded.

\subsection*{62. General localizers and the failure of the first characterization (pp.~156--159)}
The new entry, visibly dated 6.3. rather than the witness's 6.4., retracts the attempt to choose two independent classes of arrows. With one saturated $\underline W$, the comparison diagram proves both $\underline W=f_i^{-1}(\underline W)$ and a canonical isomorphism of the induced $\overline f_i$ with the identity on $\underline W^{-1}(Cat)$; the middle arrow in the resulting chain is used through its inverse after localization. The hypotheses combine relative asphericity of $A$, contractibility of $i(a)$, the interval condition, the fibered-category condition, and the slice criterion developed in section 61. They require no restrictive 0-admissibility assumption, so even the class of all arrows satisfies them. This forces the main conceptual revision: the argument supplies a general modelizing theory rather than a unique characterization of ordinary weak equivalences. Grothendieck sketches coefficient-dependent, twisted and truncated cohomological variants, adopts $H_{\underline W}$ or $(Hot)_{\underline W}$ for the localized category, and asks whether ordinary weak equivalences form the smallest acceptable class, so that other examples would be localizations of the usual homotopy category. He also entertains finer notions, including simple homotopy types, as possibilities rather than results. These reflections motivate section 63's terminology and the later work with fundamental localizers. The marginal note on p.~159 is consequential: it explicitly says that the impression is incorrect and suggests instead the biggest class under an extra requirement; the two words crossed by body typing remain doubtful in the checked text. The witness omits this qualification, as well as the inserted possibly before refinements, and therefore should not be used to turn the exploratory paragraph into a settled assertion. Three smaller source slips remain visible on p.~157: the definition prints $i_Wi^*$ instead of the previously defined $i_Ai^*$, the three-map diagram is called (12) instead of (11), and saturation is called condition b) although it was a). The lower comparison row has primed objects, which the supplied witness row loses. The guide identifies these dependencies without changing the source formulas or replacing the author's subsequent change of mind with a modern theorem.


\subsection*{63. Relative asphericity and terminology under revision (pp.~159--163)}
Grothendieck revisits terminology after allowing general classes $\underline W$. In a category $M$ with a final object $e$, an object with $X\to e$ in $\underline W$ is renamed a $\underline W$-object, and $M(\underline W)$ denotes the full subcategory of such objects; a map is $\underline W$-aspheric when every base change belongs to $\underline W$. The change responds to a real ambiguity: in a topos, relative weak equivalence to its final object need not make the slice topos absolutely aspheric unless the original topos is itself aspheric. For categories the induced morphism of presheaf topoi leads to the weaker condition that every $A_{/b}\to B_{/b}$ belongs to $\underline W$, called weak $\underline W$-asphericity; it does not require arbitrary base change in $(Cat)$. Page 162 compares these notions through the category-of-elements functor: under the slice criterion f), a map $u:F\to G$ of presheaves is universally in $W_A$ exactly when $i_A(u)$ is weakly $\underline W$-aspheric. This makes section 61's frequently used criterion into an explicit standing hypothesis and prepares section 64's localization axiom. The section ends by distinguishing the interval homotopy structure $h_W$ from the canonical $h_M$ of p.~123: their equality is expected in the intended examples, but the author expressly allows a setting in which additional assumptions would be unnecessary and the chosen structure must instead be specified. The terminological decision is itself qualified in the margin of p.~160 and again on p.~164, where the author keeps saying $\underline W$-aspheric for a category whose map to $e$ lies in $\underline W$. The witness's $G'=A\to G$ on p.~162 must be read against the scan's lowercase representable $a$, and its $A_F\to A_G$ omits the visible slice slashes. The same page prints the class $W_{(Cat)}$, not merely $(Cat)$. At p.~163 the overwritten verb is have, not the earlier transcription's doubtful had; the homotopy-structure discussion is not a retraction of the earlier definition. These distinctions explain the subsequent axioms without imposing a uniform modern meaning on every occurrence of aspheric.

\subsection*{64. Basic localizers, optional axioms and a degeneracy claim\newline (pp.~163--168)}
This section fixes the general framework prepared in sections 61--63. A saturated class satisfies identities, two-out-of-three and the two-sided composite condition; strong saturation means that it is exactly the inverse image of the isomorphisms of its localization, and is sought as a consequence rather than imposed initially. The provisional list consists of L1 saturation, L2 interval homotopy invariance, L3 the final-object condition, its special case L3-prime for the interval, L4 the slice localization criterion and L5 the fibered-category criterion. A basic localizer is then defined by L1, L3, L4, which supply the homotopy axiom and recover the previous test-category theory, initially with final objects in the values of test functors. Dropping that final-object restriction motivates L5 and the next section's promised comparison theorems. The oblique note on p.~165 qualifies the implication from L5 to localization: it invokes $W=W^o$, and otherwise asks for L5 plus a dual axiom concerning cofibrations. This is materially different from the earlier transcription's supposed reference to p.~30, and the witness omits the qualification. Grothendieck next considers optional requirements: 0-admissibility and its weaker connectedness condition, preservation of sums and products by localization, invariance under passage to opposite categories, and a stronger relative slice criterion L4-prime. A French note on p.~167 says that the duality assertion is true for fundamental localizers and refers to the duals of L5 and L4-prime; the overlapped words and final label remain doubtful in the checked source. The section concludes by extending all test terminology relative to a fixed $\underline W$, while reserving $(Hot)$ for ordinary homotopy types and using $(Hot)_{\underline W}$ in general. A substantive gap occurs in the proposition on p.~166: equality of all parallel arrows after localization does not imply that the localization is terminal. Let $q:(Cat)\to[1]$ record whether a category is empty, and let $W=q^{-1}(\operatorname{Iso}[1])$. This class consists of all functors between nonempty categories together with the identity of the empty category, is saturated, and satisfies L2 because product with the nonempty interval preserves emptiness. The natural maps from a nonempty category to the one-object category exhibit $W^{-1}(Cat)\simeq[1]$, with distinct empty and nonempty objects. The discrete two-point category is universally $W$-aspheric and makes any parallel pair homotopic for $h_W$, whereas its two points are not connected by the canonical homotopy relation. Thus $h_W\ne h_{(Cat)}$ although the localization is not terminal. The source argument correctly collapses parallel arrows but overlooks the remaining empty-to-nonempty distinction; the readers retain the statement. Smaller cautions are p.~163's composition $gf$, reversed to $fg$ by the witness; the source's basic modelizer in the propositions on pp.~166--167, where the class under discussion is a basic localizer; and p.~167's literal for any $A$ in $C$, where the product criterion requires arbitrary $A$ in $(Cat)$. The clearly typed so si on that page alone is corrected as an explicitly recorded letter transposition. These limitations belong alongside the evolving hypotheses, before the cleaner test-functor review beginning on p.~168.


\subsection*{65. Model-preserving functors and the general test-category theorem\newline (pp.~168--179)}
After the long investigation of homotopy structures, Grothendieck returns to the comparison problem announced around p.~88. A model-preserving functor both reflects and preserves the chosen weak equivalences and induces an equivalence of localizations; neither category is assumed beforehand to be a modelizer. With a fixed basic localizer $\underline W$, he reviews the category-of-elements adjunction $i_A\dashv i_A^*$, the class $\underline W_A=i_A^{-1}(\underline W)$, and the difference between absolute asphericity of $A_{/F}$ and universal asphericity of the projection $F\to e_{\widehat A}$. The latter is tested on products with representables, while general universally weakly equivalent maps are tested on representable base changes. Part A characterizes closure of aspheric presheaves under products; the handwritten addition in Definition 1 also requires $A$ itself to be aspheric before calling its presheaf topos totally aspheric. This is a substantive added hypothesis, retained in the readers. The subsequent discussion explains how contraction of representables and asphericity of interval products can verify the criterion, and proposes the Mayer--Vietoris condition L e) for a cover of a category by two full subcategories closed under passage to sources. The marginal note calls these open subobjects and identifies their presheaf topoi as open subtopoi, not merely open sets. Whether the extra axiom follows from the preceding axioms remains a question here. Parts B--D characterize weak, local and ordinary test categories through the adjunction, the Lawvere contractor and products with an interval; the right marginal addition on p.~173 records a further local criterion, preservation by $i_A^*$ of weak equivalences as universally $\underline W_A$-aspheric maps, and its restricted final-object version. On p.~174 Grothendieck explicitly notices that a hypothesis on $i_A^*$ may have been omitted in the theorem on p.~46; the guide records this reconsideration without silently adding the hypothesis to the earlier text. Part E extends weak and ordinary test functors to $i:A\to M$ for a category endowed with a saturated class $W$: the three basic requirements are model preservation by $i^*$, the weak-test property of $A$, and its asphericity. The induced composite $f_i=i_Ai^*$ provides an equivalent comparison. Theorem 1 then specializes to $M=(Cat)$, assumes L5 and contractibility of every $i(a)$, and characterizes the test property by universal asphericity of $i^*(\Delta_1)$ together with asphericity of $A$. The proof uses nonemptiness to separate the interval endpoints and the key lemma on p.~156 for model preservation; it does not need a final object in every $i(a)$. The intervening contractibility criterion is weaker than the false final-object converse on p.~151 discussed in section 60. The local variant and corollary on p.~178 separate the interval condition from the additional asphericity of $A$; a boxed marginal claim about a model-preserving functor from $(Cat)_{/A}$ is explicitly marked to check and remains provisional. Part F assumes total asphericity and proves Theorem 2: the test property is equivalent successively to model preservation, preservation of weak equivalences, preservation of aspheric objects, preservation of contractible objects, and the single test on $\Delta_1$. The revised implication diagram and the generating-contractor argument give the equivalence of the last two tests; universal asphericity and contractibility are related using the total-asphericity assumption. This completes the general comparison whose significance section 66 reassesses. Several literal formulas need care: p.~169 omits the prime on the target localization and later prints $\underline W_{A^\wedge}$ where the defined class is $\underline W_A$; its displayed target $A_{/G}$ is a typed correction over $A_{/F}$ that the witness misses. On p.~172 the incomplete correction of the to a symbol with subscript $A$, and the doubtful nearby two-$A$ formula, are not reconstructed from mathematical expectations. In the remark on p.~177, condition a) literally repeats $(i_{/a})^*:M\to(A_{/a})^\wedge$, also used in b), although the explanation that $A$ is a local test category calls for $i_{A_{/a}}^*:(Cat)\to(A_{/a})^\wedge$; the source expression is preserved and this likely index/domain slip is confined to the guide. Page 178 omits the subscript 1 from a displayed $\Delta$, and p.~179's condition (iv) literally says $\underline W_A$-objects, with its explanatory parenthesis supplying the asphericity meaning. The strict-total-asphericity qualification on p.~173 should also be read with section 64's caution about the degeneracy argument. These distinctions preserve the developing hypotheses and the author's corrections without converting the review into a retrospectively uniform theorem statement.
